\subsection{Anatomy of Scope Failures}

Scope failures reveal how the compiler classified names before runtime. Each error type corresponds to a specific mismatch between programmer expectation and static binding rules.

\subsubsection{Runtime Execution Traces and Bytecode Analysis of \texttt{UnboundLocalError}}

If a name is assigned anywhere in a function, the compiler marks it local for the entire function. A read before that assignment executes \texttt{LOAD\_FAST} on an uninitialized slot.

%<<<PROGRAM:programs/ana01>>>


Bytecode uses \texttt{LOAD\_FAST} for \texttt{value}, not \texttt{LOAD\_GLOBAL}, because the assignment later in the body forced local classification.

\subsubsection{Analyzing the Mechanics of Conflicting Local and Global Names}

The same name cannot serve as both an unread global and a written local without \texttt{global} declaration. The compiler chooses local when any assignment exists; global reads of that name then fail unless \texttt{global name} is declared.

\subsubsection{Name Resolution Failure: \texttt{NameError} vs. \texttt{UnboundLocalError}}

\texttt{NameError}: no binding exists in any searched namespace. \texttt{UnboundLocalError}: a local slot exists but has not been assigned yet. Both occur during name resolution, before the intended operation runs.
